
\section{}
\setcounter{subsection}{1}
\subsection{Marco de Referencia}
\subsubsection{Base Legal-Técnica}
El diagnóstico de seguridad se fundamenta en dos marcos normativos clave:
\begin{itemize}
	\item \textbf{LGPDPSO (Art. 6, 10, 17, 24-36)}: Establece principios y medidas técnicas obligatorias para protección de datos personales en instituciones públicas \cite{lgbdpsopdf}
	\item \textbf{NIST CSF 2.0}: Marco de ciberseguridad con funciones básicas para gestión de riesgos tecnológicos \cite{nist2023}
\end{itemize}

\subsubsection{Modelo de Evaluación NIST CSF 2.0}
Aplicación práctica del marco mediante mapeo de controles:
\begin{table}[h]
	\centering
	\caption{Correlación NIST CSF 2.0 - Hallazgos en CIITA}
	\begin{tabular}{|l|l|l|}
		\hline
		\textbf{Función NIST} & \textbf{Requerimiento} & \textbf{Estado en CIITA} \\ \hline
		Identificar & Inventario de activos & Incompleto (Art. 27 LGPDPSO) \\ 
		Proteger & Cifrado de datos & No implementado (Art. 25 LGPDPSO) \\ 
		Detectar & Monitoreo continuo & Sistema inoperante \\ 
		Responder & Plan de incidentes & No documentado (Art. 33 LGPDPSO) \\ 
		Recuperar & Copias de seguridad & Manuales/no verificadas \\ \hline
	\end{tabular}
\end{table}

\subsubsection{Requisitos Técnicos de LGPDPSO}
Principales obligaciones técnicas según los artículos relevantes:
\begin{enumerate}
	\item \textbf{Artículo 10}: Implementar principios de licitud, consentimiento y finalidad en tratamiento de datos
	\item \textbf{Artículo 25}: Establecer medidas de seguridad administrativas, físicas y técnicas
	\item \textbf{Artículo 27}: Mantener inventario actualizado de sistemas de tratamiento de datos
	\item \textbf{Artículo 29}: Elaborar documento de seguridad con análisis de riesgos
	\item \textbf{Artículo 34}: Notificar brechas de seguridad en \textless72 horas
\end{enumerate}

\textbf{Incumplimientos críticos en CIITA:}
\begin{itemize}
	\item Falta de inventario de datos personales (Art. 27.III)
	\item Falta de conocimiento de confidencialidad básica en empleados (Art. 24.III y 36)
	\item Falta de mecanismos básicos para proteger la información (Art. 25 y 26)
\end{itemize}

\subsubsection{Contexto Operativo del CIITA}
Como único centro del \textbf{IPN} especializado en integración tecnológica en la frontera norte, el CIITA:
\begin{itemize}
	\item Capacita a \textbf{350+ estudiantes anuales} en automatización industrial.
	\item Colabora con \textbf{15+ maquiladoras} en proyectos de I+D en sectores aeroespacial y médico.
	\item Requiere \textbf{2.5 millones de MXN} para modernizar infraestructura (18\% del presupuesto anual).\cite{ipn2023}
\end{itemize}

El CIITA debe cumplir adicionalmente con:
\begin{itemize}
	\item \textbf{Ley de Firma Electrónica Avanzada}: Para validar integridad de comunicaciones oficiales
	\item \textbf{Ley General de Transparencia}: En lo relativo a protección de datos públicos no sensibles
\end{itemize}

\subsubsection{Riesgos Técnicos Identificados}
\begin{itemize}
	\item \textbf{Red no segmentada}: Dispositivos en la red de invitados acceden a recursos internos sin restricciones.
	\item \textbf{Equipos obsoletos}: 30+ computadoras con sistemas operativos sin soporte (ej: Windows 10) en laboratorios.
	\item \textbf{Gestión de TI remota}: Soporte técnico desde Chihuahua causa retrasos en resolver incidentes críticos.
\end{itemize}

\subsubsection{Impacto en Certificaciones}
La infraestructura actual del CIITA contradice requisitos de:
\begin{itemize}
	\item \textbf{ISO 9001}: Falta de controles para garantizar calidad operativa.
	\item \textbf{ISO 27000}: Configuraciones de red inseguras y datos expuestos.
\end{itemize}

\subsubsection{Metodología de Análisis}
El diagnóstico se realizó mediante:
\begin{itemize}
	\item \textbf{Observación de tráfico}: Identificación de dispositivos no autorizados en subredes críticas mediante herramientas de escaneo pasivo (ej: Wireshark).
	\item \textbf{Pruebas de conectividad}: Verificación de políticas de firewall mediante intentos de acceso a sitios web comúnmente bloqueados (ej: redes P2P, contenido malicioso).
	\item \textbf{Revisión documental}: Análisis de políticas de seguridad existentes proporcionadas por el personal del CIITA.
	\item \textbf{Entrevistas}: Consultas al personal técnico para validar hallazgos y contextualizar limitaciones operativas.
\end{itemize}